\documentclass[10pt,twocolumn,letterpaper]{article}

\usepackage[utf8]{inputenc}
\usepackage[T1]{fontenc}
\usepackage{amsmath}
\usepackage{amssymb}
\usepackage{graphicx}
\usepackage{xcolor}
\usepackage{hyperref}
\usepackage{cite}
\usepackage{array}
\usepackage{booktabs}
\usepackage{fancyhdr}
\usepackage{geometry}
\usepackage{tabularx}

\geometry{margin=0.75in}

\title{\textbf{Self-Supervised Monocular Depth Estimation from Videos: A Review}}

\author{
\begin{minipage}[t]{0.45\textwidth}
\centering
\textbf{Shantanu Shukla}\\
Centre for Development of Advanced Computing (C-DAC)\\
Noida, India\\
\href{mailto:shantanu.shukla10@gmail.com}{shantanu.shukla10@gmail.com}
\end{minipage}
\hfill
\begin{minipage}[t]{0.45\textwidth}
\centering
\textbf{Rosy Verma}\\
Centre for Development of Advanced Computing (C-DAC)\\
Noida, India\\
\href{mailto:rosyverma@cdac.in}{rosyverma@cdac.in}
\end{minipage}
}

\date{}

\begin{document}

\maketitle

\begin{abstract}
Self-supervised monocular depth estimation from videos has quickly grown into an active research area because it promises depth perception without expensive sensors or manual labels. Instead of relying on ground-truth depth, these methods learn from the relationships between frames in ordinary video sequences and use image reconstruction as a supervisory signal. Over the last few years, this idea has produced a diverse set of approaches that differ in how they handle dynamic scenes, enforce temporal consistency, preserve fine details, and trade off accuracy against efficiency.

This paper reviews ten representative works published between 2018 and 2025 that have shaped this field. The survey starts from early methods that explicitly model object motion, then moves to architectures that simplify training objectives, and finally covers recent designs that target edge deployment or long video sequences. For clarity, the discussion is organized into three broad phases: foundational methods, techniques that refine temporal and geometric consistency, and specialized models designed for particular deployment requirements. Along the way, recurring open problems emerge, including scale ambiguity, limited generalization across domains, and the difficulty of achieving temporal stability without incurring high computational cost.

By tracing how these ideas evolved across several generations of models, this review aims to give readers a coherent picture of where self-supervised video depth estimation currently stands and what challenges remain. The paper also highlights directions that appear promising for future work, such as combining motion-aware reasoning with lightweight architectures and integrating additional sensor modalities for robust metric depth.
\end{abstract}

\noindent \textbf{Index Terms:} Self-supervised learning, monocular depth estimation, video depth, temporal consistency, deep learning

\section{Introduction}
\label{sec:intro}

Estimating depth from a single camera has long been a central challenge in computer vision. Applications such as autonomous driving, robotics, augmented reality, and surveillance all benefit from reliable depth information, yet many real-world systems cannot afford stereo rigs or active sensors such as LiDAR \cite{hartley2003multiple}. Monocular setups are cheaper and easier to deploy, but recovering 3D structure from a single image or video stream is inherently ambiguous.

Self-supervised monocular depth estimation addresses this challenge by learning directly from videos instead of relying on labeled depth maps \cite{lecun2015deep}. The key idea is that when a camera moves through a scene, the way objects shift on the image plane depends on their distance and the camera motion. By training a network to reconstruct one frame from another using a predicted depth map and estimated pose, one can indirectly supervise the depth prediction process \cite{jaderberg2015spatial,godard2019digging}. This approach removes the need for explicit depth annotations and opens the door to using large amounts of unlabeled video data.

In only a few years, self-supervised depth from videos has evolved from proof-of-concept demonstrations into a fairly mature research area with many branches. Some methods focus on making the underlying geometry more consistent by modelling object motion or separating ego-motion from dynamic elements \cite{casser2019depth}. Others prioritize temporal stability, aiming to reduce frame-to-frame flickering that can be problematic in downstream applications \cite{zhang2019exploiting,li2021enforcing}. More recent work has pushed the field toward practical deployment, designing architectures that are lightweight enough for embedded devices \cite{zhang2023lite} or robust enough to handle long video sequences \cite{chen2025video}.

Because of this diversity, it is not obvious how these contributions relate to one another or where the main gaps now lie. The goal of this paper is to bring these pieces together. We review ten influential works from 2018 to 2025 that span dynamic-scene modeling, temporal consistency, detail-preserving architectures, lightweight designs, and long-sequence processing. Rather than simply listing methods, the review emphasizes how ideas developed over time and how each method responds to limitations of earlier work.

To make the discussion easier to follow, we group the selected papers into three stages. The first stage includes foundational methods that established the feasibility of self-supervised monocular depth estimation from video and introduced explicit treatment of moving objects. The second stage looks at methods that strengthen temporal and geometric consistency through recurrent architectures, attention mechanisms, or optimization-based refinements. The final stage focuses on specialized approaches that optimize particular aspects such as efficiency, dynamic-scene robustness, or scalability to long videos. We conclude by summarizing common limitations and outlining possible directions for future research.

\section{Foundational Approaches: Handling Dynamic Scenes}
\label{sec:foundational}

\subsection{Struct2Depth (2018)}

Early self-supervised depth methods typically assumed that the world is static and that all motion in the image is caused by camera movement. In practice, this assumption is frequently violated: pedestrians walk, vehicles move independently, and objects may deform. Struct2Depth was one of the first methods to explicitly address this issue in the context of monocular video depth estimation \cite{casser2019depth}.

The central idea in Struct2Depth is to jointly predict depth, camera pose, and masks that segment independently moving objects. By decomposing motion into background and foreground components, the method avoids explaining away object motion as erroneous depth or camera pose. This is achieved by introducing instance-level motion fields alongside the global ego-motion model and using them to warp frames more accurately. As a result, the model can handle scenes with moving cars and people better than purely static assumptions would allow.

However, this capability comes at a cost. Struct2Depth relies on additional components such as pre-trained instance segmentation, which increases both complexity and computational requirements. The framework works best on short clips and controlled datasets, and its reliance on high-quality segmentation is a potential bottleneck. Nevertheless, the method clearly showed that dynamic scenes are not a corner case but a central concern that needs to be built into self-supervised depth systems.

\subsection{Monodepth2 (2019)}

Monodepth2 took a rather different path. Instead of adding extra networks to explicitly model object motion, it focused on improving the learning signal within a relatively simple architecture \cite{godard2019digging}. Two tricks, in particular, proved highly effective and have since been widely adopted.

The first is the \textit{minimum reprojection loss}. When there are multiple candidate source frames that could be used to reconstruct a target frame, Monodepth2 computes the photometric error for each candidate and, for each pixel, takes the minimum. This reduces the influence of occlusions and out-of-view regions, because only the best matching view contributes to the loss for each pixel. The second idea is \textit{auto-masking}, which identifies pixels that do not provide meaningful training signal—typically those with little or no relative motion across frames—and excludes them from the loss. Together, these mechanisms help the model focus on informative regions and reduce artifacts related to stationary or ambiguous pixels.

Monodepth2 is appealing because these ideas are conceptually simple, easy to implement, and do not significantly increase inference-time complexity. It quickly became a standard baseline for self-supervised monocular depth estimation and remains a common starting point for new methods. At the same time, its predictions are computed frame by frame, without any explicit temporal smoothing, which leads to visible flickering in video sequences. This lack of temporal consistency and the limited treatment of moving objects motivate several of the later methods in this review.

\subsection{PackNet-SfM (2019)}

While Monodepth2 emphasized training losses and simplicity, PackNet-SfM looked at the architecture itself and asked how much spatial detail can be preserved in self-supervised depth maps \cite{guizilini20203d}. The authors proposed a novel encoder–decoder design based on ``3D packing'' and ``unpacking'' operations, which reorganize feature maps in a way that preserves more fine-grained structure compared to standard downsampling.

Intuitively, the packing operation re-arranges spatial information into channel dimensions, allowing the network to process high-resolution details without increasing the spatial size of feature maps. The corresponding unpacking step then recovers a dense, high-resolution depth map. This design proved effective for generating sharp depth boundaries and preserving small objects that might otherwise be smoothed away. PackNet also explored optional velocity supervision, where ground-truth vehicle velocities are used to recover metric scale in the depth predictions rather than only relative depth.

The main downside of PackNet is its size and computational footprint. The architecture contains on the order of hundreds of millions of parameters, and inference is comparatively slow. While this is acceptable for offline processing and some research scenarios, it is not ideal for resource-constrained platforms. Still, PackNet demonstrates that architectural choices can significantly affect the level of detail and sharpness in self-supervised depth outputs.

\section{Temporal Consistency and Geometric Refinement}
\label{sec:temporal}

\subsection{Exploiting Temporal Consistency (2019)}

Monodepth2 and similar models process each frame independently at test time, which often leads to depth maps that flicker across consecutive frames even when the scene is static. Zhang et al.\ addressed this issue by explicitly incorporating temporal memory into the depth network \cite{zhang2019exploiting}. Their model introduces a spatio-temporal recurrent structure that uses Long Short-Term Memory (LSTM) units to propagate information over time.

By reusing temporal context, the network is encouraged to produce more stable depth maps from frame to frame. The authors also proposed metrics to quantify temporal consistency, making it possible to evaluate not just per-frame accuracy but also temporal smoothness. In practice, their model delivers real-time video depth while visibly reducing temporal fluctuations compared to single-frame baselines.

The trade-off is that recurrent models are more complex and can struggle with very long sequences. LSTMs have a limited effective memory window, and as sequences get longer, distant past information tends to be forgotten or overshadowed. Additionally, the inclusion of recurrent units and adversarial losses increases training and inference cost compared to lighter baselines.

\subsection{Self-Attention and Discrete Disparity Volume (2020)}

Another line of work explored using self-attention mechanisms to capture long-range dependencies in images. Johnston and Carneiro integrated self-attention into a self-supervised monocular depth framework, along with a discrete disparity volume representation \cite{johnston2020self}. The key intuition is that depth estimation benefits from global context: a pixel on a textureless wall, for example, might be better interpreted if the network can attend to edges or objects elsewhere in the image.

In their method, features are processed through self-attention blocks that allow each location to weigh contributions from all other locations, providing a richer representation of scene structure. At the same time, depth is represented as a discrete probability distribution over disparity levels rather than a single predicted value. This makes it possible to express uncertainty, which is especially helpful around occlusion boundaries and regions with weak texture.

These additions improve depth quality and robustness, but they mainly operate within each frame. While the method lays groundwork for using attention and uncertainty in self-supervised depth, it does not yet directly address the temporal dimension of video sequences. This gap is taken up by methods that combine such ideas with temporal constraints or multi-frame processing.

\subsection{Consistent Video Depth via Geometric Optimization (2020)}

Instead of relying solely on learned networks, Luo et al.\ asked whether classical multi-view geometry could be used to refine depth predictions and enforce consistency \cite{luo2020consistent}. Their approach first uses a neural network to produce an initial depth estimate, then applies a geometric optimization step at test time to improve temporal and spatial coherence.

This optimization enforces consistency between depth, pose, and optical flow across frames by penalizing violations of epipolar geometry and other geometric relationships. The result is a sequence of depth maps that better respect the underlying 3D structure of the scene. From a conceptual standpoint, this work is notable because it demonstrates that learned predictions and geometric constraints can complement each other rather than compete.

The downside is computational cost. The optimization procedure is relatively expensive and has to be run at test time, making it unsuitable for real-time applications or long video sequences where many frames need to be processed. Nevertheless, the method provides a useful reference point for what can be achieved when geometry is enforced explicitly.

\subsection{Training-Time Temporal Consistency (2021)}

Li et al.\ proposed a more lightweight alternative for improving temporal consistency that does not require recurrent networks or test-time optimization \cite{li2021enforcing}. Their idea is to incorporate a temporal consistency loss during training, using optical flow to relate depth predictions across frames.

The temporal loss encourages depth at a pixel in one frame to be consistent with the depth of the corresponding pixel in a neighboring frame, according to the estimated optical flow, unless motion suggests otherwise. If two frames show the same part of the scene and the optical flow is reliable, their depths should agree up to expected changes. This encourages smoother temporal behavior without changing the test-time architecture, which can remain as simple as Monodepth2.

In practice, this approach improves temporal stability at little or no additional inference cost. The main sensitivity lies in the quality of the optical flow: if flow estimates are inaccurate, the temporal constraint can become unreliable. Even so, this line of work demonstrates that a lot of temporal consistency can be baked in at training time, without making test-time models significantly heavier.

\section{Specialized Architectures and Deployment-Oriented Methods}
\label{sec:specialized}

\subsection{Lite-Mono: Lightweight Depth for Edge Devices (2023)}

As self-supervised depth models became more accurate, another question grew more pressing: how can these models be deployed on resource-constrained hardware such as drones, mobile phones, or embedded systems? Lite-Mono addresses this practical concern by designing a depth network that is both compact and effective \cite{zhang2023lite}.

Lite-Mono combines convolutional backbones with efficient attention-like modules to capture multi-scale context without incurring the cost of full self-attention. The architecture is carefully constructed to keep parameter count and FLOPs low, while still leveraging ideas from modern vision transformers and multi-scale feature extraction. In experiments, Lite-Mono achieves competitive or better accuracy compared with much larger models, while using only a fraction of the parameters and running in real time on edge hardware.

One limitation is that Lite-Mono focuses primarily on single-frame depth estimation and does not explicitly incorporate temporal consistency or motion modeling. As such, it is an excellent building block for efficient depth estimation, but additional components would be required to handle dynamic scenes or to stabilize predictions across video sequences. Still, it serves as an important step toward making self-supervised depth practical beyond desktop GPUs.

\subsection{ManyDepth2: Motion-Aware Depth in Dynamic Scenes (2024)}

Dynamic scenes remain a major challenge for monocular depth estimation. ManyDepth2 targets this problem directly by designing a motion-aware architecture that can disentangle camera motion from object motion and use this information to improve depth estimates in dynamic environments \cite{zhou2024manydepth2}.

At a high level, ManyDepth2 constructs several different ``views'' of the scene for each pixel. One view assumes a static world with only camera motion, another uses a motion-compensated reference that attempts to remove object motion, and a third leverages optical flow-based warping. These perspectives are aggregated through learned attention weights that decide, for each pixel, which perspective to trust most. For example, static background regions may rely predominantly on the standard geometric view, while pixels on moving vehicles may lean more on the motion-compensated or flow-based perspectives.

This multi-perspective fusion leads to significantly better performance on moving objects and more accurate camera pose estimation in dynamic scenes. The trade-off is increased complexity and computational cost: ManyDepth2 is heavier than baseline methods and is not yet suited for very low-power devices. Nonetheless, it shows how explicitly reasoning about motion can substantially improve depth estimation in the kinds of scenes that matter in practice, such as traffic and crowded urban environments.

\subsection{Video Depth Anything: Long-Sequence Depth (2025)}

Most earlier methods focus on short clips or per-frame evaluation, but many real-world applications require processing long video streams—potentially minutes or hours of footage. Video Depth Anything (VDA) tackles this by extending self-supervised depth methods to operate efficiently over long sequences while maintaining temporal consistency \cite{chen2025video}.

Instead of treating each frame independently, VDA uses a combination of key-frame selection, temporal feature propagation, and temporal regularization. Key frames are processed more thoroughly, while intermediate frames reuse information propagated from these key frames, which keeps computation manageable. Temporal gradient losses or similar regularizers are applied to discourage abrupt changes in depth over time, especially in regions where the scene structure is stable.

This design allows VDA to handle super-long videos in a way that balances computational cost, spatial accuracy, and temporal smoothness. The approach still tends to rely on reasonably powerful hardware, and key-frame strategies can introduce mild artifacts at segment boundaries if not handled carefully. Even so, VDA represents an important move toward depth estimation systems that can operate on realistic, long-duration video streams rather than only short snippets.

\section{Comparative Evolution and Synthesized Insights}
\label{sec:comparative}

\subsection{Evolution Across Three Phases}

Although the ten methods reviewed above differ in many ways, they can be loosely grouped into three phases that reflect how the field has evolved over time. Table~\ref{tab:evolution} summarizes this progression.

\begin{table*}[t]
\centering
\footnotesize
\begin{tabularx}{\textwidth}{|l|c|p{3cm}|X|}
\hline
\textbf{Phase} & \textbf{Years} & \textbf{Primary Focus} & \textbf{Representative Works} \\
\hline
Foundational & 2018--2019 & Feasibility and handling of dynamic scenes & Struct2Depth \cite{casser2019depth}, Monodepth2 \cite{godard2019digging}, PackNet-SfM \cite{guizilini20203d} \\
\hline
Refinement & 2019--2021 & Temporal stability and geometric consistency & Zhang et al.\ \cite{zhang2019exploiting}, Johnston \& Carneiro \cite{johnston2020self}, Luo et al.\ \cite{luo2020consistent}, Li et al.\ \cite{li2021enforcing} \\
\hline
Specialization & 2023--2025 & Efficiency, dynamic-scene robustness, long videos & Lite-Mono \cite{zhang2023lite}, ManyDepth2 \cite{zhou2024manydepth2}, Video Depth Anything \cite{chen2025video} \\
\hline
\end{tabularx}
\caption{High-level evolution of self-supervised monocular depth estimation from videos, grouped into three phases from early feasibility to deployment-oriented specialization.}
\label{tab:evolution}
\end{table*}

The foundational phase established that self-supervised learning from monocular videos is viable and that dynamic scenes require explicit attention. The refinement phase recognized that temporal consistency and geometric correctness are key to making these models useful in video settings. The specialization phase then pushed the field in several directions at once: toward smaller models, toward better handling of complex motion, and toward scalability to longer sequences.

\subsection{Key Conceptual Lessons}

Looking across these works, several recurring themes become apparent:

\begin{itemize}
    \item \textbf{Dynamic scenes cannot be ignored.} Methods like Struct2Depth and ManyDepth2 show that independently moving objects, such as vehicles and pedestrians, can significantly degrade depth quality if treated as static background. Explicit motion modeling or multi-perspective fusion helps address this.

    \item \textbf{Loss design matters as much as architecture.} Monodepth2 demonstrates how relatively simple changes to the training objective, such as minimum reprojection loss and auto-masking, can yield large gains without making the model much more complex.

    \item \textbf{Temporal consistency is a separate axis of quality.} Per-frame accuracy does not guarantee visually stable video depth. Methods that introduce recurrent structures, training-time temporal losses, or geometric consistency constraints highlight the importance of explicitly considering time.

    \item \textbf{Detail preservation and uncertainty are important.} Architectural choices like those in PackNet and the use of discrete disparity volumes with uncertainty estimation show that sharpness and reliable confidence estimates are valuable aspects of depth prediction.

    \item \textbf{Efficiency and scalability drive practical impact.} Lite-Mono and Video Depth Anything illustrate that, beyond accuracy, the community is increasingly concerned with making models run on real hardware and on realistic, long-duration video streams.
\end{itemize}

These themes suggest that no single method solves all aspects of self-supervised video depth. Instead, each contribution occupies a different point in the design space, and there is room for more integrated approaches that combine these strengths.

\section{Open Challenges and Limitations}
\label{sec:challenges}

Despite rapid progress, several fundamental challenges remain open.

\subsection{Metric Scale Ambiguity}

Self-supervised monocular depth estimation typically recovers depth only up to an unknown scale factor. Without additional information, the network cannot distinguish between a small object that is nearby and a large object that is farther away if their projections look similar. Some methods partially address this via velocity supervision or scale alignment during evaluation, but a general, robust solution often requires extra sensors (such as IMU or GPS) or sparse depth measurements. Integrating such cues in a principled way remains an active area of research.

\subsection{Domain Generalization}

Many methods are developed and tested on automotive datasets such as KITTI or similar urban driving scenarios. When applied to very different environments—indoor scenes, night-time imagery, aerial videos, or underwater footage—performance often drops considerably. Bridging this domain gap, whether by better data augmentation, domain adaptation techniques, or larger-scale pretraining, is crucial for making self-supervised depth estimation broadly usable.

\subsection{Balancing Efficiency and Dynamic-Scene Performance}

Models optimized for efficiency, such as Lite-Mono, tend to focus on single-frame inference and may underperform in highly dynamic environments. Conversely, motion-aware models like ManyDepth2 are more accurate on moving objects but are heavier and more expensive to run. Designing architectures that are both lightweight and robust to complex scene motion is still challenging and likely requires new ideas in how motion, depth, and features are represented.

\subsection{Dependence on Optical Flow Quality}

Several approaches use optical flow either as an input signal or to define temporal consistency losses. When the flow is accurate, it provides valuable information about motion and correspondence. However, flow estimation itself is a difficult problem, especially in low-texture regions, near occlusion boundaries, or under fast motion. Errors in flow can propagate into depth estimation and consistency losses, making the system brittle.

\subsection{Long-Sequence Consistency}

Methods like Video Depth Anything show that processing long sequences is feasible, but maintaining consistent depth across thousands of frames remains challenging. Key-frame strategies and temporal regularization help, yet they can introduce subtle discontinuities at segment boundaries. More robust mechanisms for global temporal coherence, possibly based on learned long-term memory or global constraints, are still needed.

\section{Future Directions}
\label{sec:future}

Looking ahead, several directions appear particularly promising.

\subsection{Unified Models that Balance Multiple Objectives}

Most current methods optimize strongly for one or two aspects, such as accuracy, temporal consistency, or efficiency. A natural next step is to design architectures and training schemes that explicitly target a balanced trade-off between these objectives. For example, combining ideas from Lite-Mono and ManyDepth2 could lead to motion-aware depth networks that are still compact enough for embedded deployment.

\subsection{Multi-Modal Sensor Fusion}

In many real-world systems, cameras are not the only sensors available. Inertial Measurement Units, GPS, wheel odometry, or sparse LiDAR all provide complementary information that can help resolve scale, improve robustness, or handle challenging lighting conditions. Incorporating these sensors in a self-supervised manner—so that depth and motion predictions remain consistent with all available data—is a promising research avenue.

\subsection{Better Domain Adaptation and Pretraining}

Large-scale pretraining on diverse video datasets, followed by fine-tuning on specific domains, has worked well in other areas of computer vision. Applying similar strategies to self-supervised depth estimation, and combining them with explicit domain adaptation techniques, may improve generalization to new environments without requiring extensive labeled data.

\subsection{Uncertainty and Reliability Estimation}

As depth predictions are used in safety-critical applications, it becomes increasingly important not just to estimate depth, but also to know when the model is uncertain. Approaches that provide well-calibrated uncertainty estimates or that can detect out-of-distribution inputs are likely to be valuable in practice.

\section{Discussion}
\label{sec:discussion}

Taken together, the methods surveyed in this paper tell a story of a field that has matured quickly while still leaving important questions open. Early work showed that it is possible to learn depth from monocular video without ground-truth labels, and that dynamic scenes require special handling. Subsequent methods refined the training objectives and architectures to improve temporal consistency, geometric validity, and detail preservation. More recent contributions have pushed toward real-world constraints such as limited compute, dynamic environments, and long video durations.

One pattern that emerges is that progress often comes from recognizing a specific weakness in existing methods and addressing it directly. For example, Monodepth2 improves robustness by carefully designing the photometric loss; Li et al.\ reduce flickering by adding a temporal consistency loss; Lite-Mono makes the model lighter without sacrificing much accuracy; ManyDepth2 explicitly confronts the challenge of moving objects. Each of these contributions shifts the trade-off surface between accuracy, stability, complexity, and practicality.

At the same time, this specialization means that no single method covers all desirable properties. There is still a gap between what is feasible in controlled benchmarks and what is needed in fully deployed systems that must operate over long durations, in varied environments, and under tight resource constraints. Bridging this gap will likely require combining insights from multiple lines of work, along with stronger integration of geometry, motion, and multi-modal sensing.

\section{Conclusion}
\label{sec:conclusion}

This paper has reviewed ten key contributions to self-supervised monocular depth estimation from videos, spanning roughly seven years of research. The methods reviewed cover dynamic-scene modeling, temporal consistency, architectural design for detail preservation, efficiency-focused models, and long-sequence processing. By organizing these works into three phases and highlighting the ideas that connect them, we have aimed to provide a clearer picture of how the field has evolved.

Despite strong progress, important challenges remain. Metric scale ambiguity, limited domain generalization, the tension between efficiency and dynamic-scene robustness, reliance on optical flow, and long-sequence consistency are all areas where more work is needed. Future research that blends motion-aware reasoning with lightweight architectures, leverages additional sensor modalities, and explicitly targets multiple objectives at once may bring self-supervised depth estimation closer to widespread real-world deployment.

Ultimately, the trajectory of the field suggests that self-supervised depth from videos will continue to be an active area of research, with growing impact in both academic and applied settings.

\begin{thebibliography}{99}

\bibitem{godard2019digging}
C.~Godard, O.~Mac~Aodha, M.~Firman, and G.~J.~Brostow, ``Digging Into Self-Supervised Monocular Depth Estimation,'' in \textit{Proc. IEEE/CVF Int. Conf. Comput. Vis. (ICCV)}, 2019, pp.~3828--3838.

\bibitem{hartley2003multiple}
R.~Hartley and A.~Zisserman, \textit{Multiple View Geometry in Computer Vision}, 2nd ed. Cambridge, U.K.: Cambridge Univ. Press, 2003.

\bibitem{lecun2015deep}
Y.~LeCun, Y.~Bengio, and G.~Hinton, ``Deep learning,'' \textit{Nature}, vol.~521, no.~7553, pp.~436--444, 2015.

\bibitem{jaderberg2015spatial}
M.~Jaderberg, K.~Simonyan, A.~Zisserman, and K.~Kavukcuoglu, ``Spatial transformer networks,'' in \textit{Proc. Int. Conf. Mach. Learn. (ICML)}, 2015, pp.~2017--2025.

\bibitem{casser2019depth}
V.~Casser, S.~Pirk, R.~Mahjourian, and A.~Angelova, ``Depth prediction without the sensors: Leveraging structure for unsupervised learning from monocular videos,'' in \textit{Proc. AAAI Conf. Artif. Intell.}, vol.~33, no.~1, 2019, pp.~8001--8008.

\bibitem{guizilini20203d}
V.~Guizilini, R.~Ambrus, S.~Pillai, A.~Raventos, and A.~Gaidon, ``3D packing for self-supervised monocular depth estimation,'' in \textit{Proc. IEEE/CVF Conf. Comput. Vis. Pattern Recognit. (CVPR)}, 2020, pp.~2485--2494.

\bibitem{zhang2019exploiting}
L.~Zhang, M.~Zhu, X.~Liu, S.~Chandrasekaran, Y.~Wang, and Y.~Shi, ``Exploiting temporal consistency for real-time video depth estimation,'' in \textit{Proc. IEEE/CVF Int. Conf. Comput. Vis. (ICCV)}, 2019, pp.~1637--1645.

\bibitem{johnston2020self}
M.~Johnston and C.~L.~Carneiro, ``Self-supervised monocular trained depth estimation using self-attention and discrete disparity volume,'' in \textit{Proc. IEEE/CVF Conf. Comput. Vis. Pattern Recognit. (CVPR)}, 2020, pp.~12088--12097.

\bibitem{luo2020consistent}
C.~Luo, Z.~Zhu, Z.~Yuille, L.~Sigal, R.~Wang, and W.~Zeng, ``Consistent video depth estimation,'' in \textit{Proc. SIGGRAPH}, 2020, pp.~52:1--52:13.

\bibitem{li2021enforcing}
Z.~Li, X.~Guo, H.~Zhu, S.~Yang, S.~Zhang, and N.~Zheng, ``Enforcing temporal consistency in video depth estimation,'' in \textit{Proc. IEEE/CVF Int. Conf. Comput. Vis. Workshops (ICCVW)}, 2021, pp.~1--10.

\bibitem{zhang2023lite}
N.~Zhang, F.~Nex, G.~Vosselman, and N.~Kerle, ``Lite-Mono: A lightweight CNN and Transformer architecture for self-supervised monocular depth estimation,'' in \textit{Proc. IEEE/CVF Conf. Comput. Vis. Pattern Recognit. (CVPR)}, 2023, pp.~16--25.

\bibitem{zhou2024manydepth2}
K.~Zhou, J.~Bian, J.~Zheng, Q.~Xie, N.~Trigoni, and A.~Markham, ``ManyDepth2: Motion-aware self-supervised monocular depth estimation in dynamic scenes,'' in \textit{Proc. IEEE/CVF Conf. Comput. Vis. Pattern Recognit. (CVPR)}, 2024, pp.~1--12.

\bibitem{chen2025video}
S.~Chen, H.~Guo, S.~Zhu, F.~Zhang, Z.~Huang, J.~Feng, and B.~Kang, ``Video Depth Anything: Consistent depth estimation for super-long videos,'' in \textit{Proc. IEEE/CVF Conf. Comput. Vis. Pattern Recognit. (CVPR)}, 2025, pp.~1--11.

\end{thebibliography}

\end{document}
